\section{Related Work}\label{sec:related}

This paper uses prior ideas on core peeling, compact clique representations, local fixed-point characterisations, parallel construction, and hierarchy output; the new ingredient is the static-\cpi{} invariant for $r=1$.

\stitle{Enumeration-based core and nucleus decomposition.}
The $k$-core was introduced by Seidman~\cite{kcoreSeidman}; Batagelj and Zaver\v{s}nik~\cite{coreAlg} gave the linear-time peeling algorithm with a bucket queue.
%
Nucleus decomposition~\cite{nucleusSariyuce14} generalises this from edges to higher-order clique witnesses by peeling $r$-cliques according to their support in $s$-cliques.
%
We specialise to the vertex-centric $r=1$ case and exploit a state invariant that does not hold for general $r$.

\stitle{Compact clique representations.}
Bron--Kerbosch search~\cite{BronKerbosch}, Tomita pivoting~\cite{tomita}, and degeneracy ordering~\cite{eppstein} are standard tools for clique enumeration; Pivoter's succinct clique tree compactly encodes clique counts~\cite{Pivoter}, and the Clique Path Index extends this line to support general nucleus decomposition with mutating hold/pivot updates during peeling~\cite{NuclearCD}.
%
The static-\cpi{} counter form of this paper shows that for vertex peeling these path edits never need to be materialised.

\stitle{Local and parallel peeling.}
$h$-index characterisations of core values were studied for local core computation~\cite{lu2016hcore,kcoreLocal,hindex}; \spin{} adapts that fixed-point view to the implicit $s$-clique witnesses encoded by the \cpi{}, and \spinStar{} is its peel-order optimisation.
%
The difference from prior local $h$-index work is not the operator but the representation: the witness multiset is held in static \cpi{} paths and updated by binomial support-loss formulas.

\stitle{Parallel construction and clique counting.}
Parallel clique search typically distributes the recursive enumeration work; \paraBuild{} follows the same seed-level decomposition.
%
Because the downstream peel never mutates path sets, construction partitions into thread-independent subproblems whose results combine by a deterministic merge --- an engineering consequence of the $r=1$ state model rather than a new clique-enumeration paradigm.

\stitle{Hierarchical cohesive subgraphs.}
Nucleus decompositions naturally form hierarchies across core levels~\cite{nucleusSariyuce14}, and elder-rule join-tree extraction is standard in computational topology~\cite{morozov2007persistence}.
%
Prior nucleus pipelines extract such hierarchies by sorting the decomposition output or revisiting clique witnesses; \buildHier{} instead reuses the preserved static \cpi{}, treating each leaf as a hyperedge born at its minimum core value, so the hierarchy is a byproduct of preserving the \cpi{} through peeling.

\stitle{Empirical graph mining methodology.}
The evaluation follows common practice in cohesive-subgraph experiments: SNAP community ground truth, large public graphs, runtime and memory comparisons, and dense synthetic stress tests~\cite{communityFinder,NuclearCD,nucleusSariyuceTODS}.
